\subsection{Proof theory for first-order logic}
\label{sec:proof-theory-fol}
\slidesrc{Proof theory (FOL): 1--4}

The source for this part is a set of university lecture notes (Math 260A, UCSD, 2012),
more technical than the rest of the course. The goal here is to understand the
\emph{ideas}: what a proof system for FOL looks like, what soundness means, and how
quantifiers are handled.

In propositional logic the simplest proof system is the truth table. In FOL there is no
analogue: the truth tables would be infinite (Section~\ref{sec:decidability}). We go
directly to \textbf{Frege-style} systems, called \textbf{Hilbert-style} systems in FOL:
\emph{many axioms, very few rules}.

\subsubsection{Frege (Hilbert-style) proofs}
\slidesrc{Proof theory (FOL): 1--2}

\begin{nbdef}{Axioms of the Frege system for FOL (page 1)}
\begin{itemize}
  \item \textbf{Propositional axioms}: all axioms of propositional Frege systems, such as
        $\varphi \limp (\psi \limp \varphi)$, where now $\varphi, \psi$ can be any FOL
        formulas.
  \item \textbf{Quantifier axioms}: for any formula $\varphi$ and any term $t$,
        \[ \varphi(t) \limp \exists x\,\varphi(x) \qquad\qquad \forall x\,\varphi(x) \limp \varphi(t) \]
        (provided $t$ is \emph{substitutable} for $x$; the notes mostly ignore this
        technical condition).
  \item \textbf{Equality axioms} (if $=$ is in the language):
        \begin{itemize}[nosep]
          \item $\forall x\,(x = x)$ \hfill(reflexivity)
          \item $\forall x \forall y\,(x = y \limp y = x)$ \hfill(symmetry)
          \item $\forall x \forall y \forall z\,(x = y \land y = z \limp x = z)$ \hfill(transitivity)
          \item $\forall \vec x\, \forall \vec y\,\big(x_1 = y_1 \land \cdots \land x_k = y_k
                \limp f(x_1, \ldots, x_k) = f(y_1, \ldots, y_k)\big)$ for every function $f$
                of arity $k$
          \item $\forall \vec x\, \forall \vec y\,\big(x_1 = y_1 \land \cdots \land x_k = y_k
                \limp (P(x_1, \ldots, x_k) \limp P(y_1, \ldots, y_k))\big)$ for every
                relation $P$ of arity $k$ \hfill(substitution of equals)
        \end{itemize}
\end{itemize}
\end{nbdef}

\begin{nbwarning}[Corrections to page 1 of the notes]
\begin{itemize}[nosep]
  \item The notes say ``$t$ should be substitutable for $t$'': it should read ``for $x$''.
  \item The last two equality axioms are printed without the hypothesis
        $x_1 = y_1 \land \cdots \land x_k = y_k$, which would make them say that $f$ gives
        the same value on \emph{any} arguments (false). The correct versions are written
        above: equal inputs give equal outputs, and equal objects have the same properties.
\end{itemize}
\end{nbwarning}

\begin{nbintuition}[Reading the quantifier axioms]
$\forall x\,\varphi(x) \limp \varphi(t)$: ``if something holds for everything, it holds for
any particular thing $t$'' (\emph{instantiation}). From ``all humans are mortal'' you may
conclude ``if Socrates is human, Socrates is mortal''.
$\varphi(t) \limp \exists x\,\varphi(x)$: ``if it holds for some particular $t$, then it
holds for something'' (\emph{generalisation to $\exists$}). From
$\mathit{Ready}(X12)$ conclude ``some monitor is ready''.
\end{nbintuition}

\begin{nbfact}{Rules of inference (pages 1--2)}
\[
\text{Modus Ponens:}\ \ \frac{\varphi \qquad \varphi \limp \psi}{\psi}
\qquad\qquad
\frac{\psi \limp \varphi}{\psi \limp \forall x\,\varphi}
\qquad\qquad
\frac{\varphi \limp \psi}{\exists x\,\varphi \limp \psi}
\]
The two quantifier rules require that \textbf{$x$ is not free in $\psi$}.
\end{nbfact}

\begin{nbintuition}[Why ``$x$ not free in $\psi$''?]
Rule 2: if $\psi \limp \varphi(x)$ has been proved and $\psi$ says nothing about $x$, then
the proof worked for an arbitrary $x$, so $\psi \limp \forall x\,\varphi$. Without the
condition we could go from $x = 0 \limp x = 0$ to $x = 0 \limp \forall x\,(x = 0)$, i.e.\
``if this number is 0 then every number is 0'': wrong. Rule 3: if $\varphi(x)$ implies
$\psi$ whatever $x$ is, then knowing that \emph{some} $x$ satisfies $\varphi$ is enough
to get $\psi$, provided $\psi$ itself does not depend on which $x$ it was.
\end{nbintuition}

\begin{nbwhiteboard}[label=wb:modus-ponens]{Modus Ponens: Socrates is mortal}{written during the FOL proof-theory lecture}
\begin{center}
\begin{tabular}{rl}
\textcircled{\scriptsize 1} & $\forall x\,.\ \mathit{human}(x) \Rightarrow \mathit{mortal}(x)$ \\
\textcircled{\scriptsize 2} & $\mathit{human}(\mathit{socrates})$ \\
\hline
 & $\mathit{mortal}(\mathit{socrates})$ \qquad \textsc{Modus Ponens}
\end{tabular}
\end{center}
\wbphoto[0.6\linewidth]{sections/first-order-logic/proof-theory-fol/assets/whiteboard-6-2.jpg}
The board shows Aristotle's syllogism (Section~\ref{sec:ai-history}) in FOL. Strictly
speaking, Modus Ponens needs an implication whose antecedent is \emph{exactly}
$\mathit{human}(\mathit{socrates})$, so a quantifier step comes first. In full, with the
Frege system above:
Writing $h$ for $\mathit{human}$, $m$ for $\mathit{mortal}$ and $s$ for
$\mathit{socrates}$:
\begin{center}
\small
\begin{tabular}{rll}
1 & $\forall x\,(h(x) \limp m(x))$ & premise \\
2 & $h(s)$ & premise \\
3 & $\forall x\,(h(x) \limp m(x)) \limp (h(s) \limp m(s))$ & quantifier axiom with $t = s$ \\
4 & $h(s) \limp m(s)$ & Modus Ponens 1, 3 \\
5 & $m(s)$ & Modus Ponens 2, 4
\end{tabular}
\end{center}
In Prolog the same knowledge is \code{mortal(X) :- human(X).} and \code{human(socrates).};
the query \code{?- mortal(socrates).} answers \code{true}: Prolog's resolution does the
instantiation ($X = \mathit{socrates}$, called \emph{unification}) and the Modus Ponens
step in one go (Section~\ref{sec:prolog}).
\end{nbwhiteboard}

\subsubsection{Provability and soundness}
\slidesrc{Proof theory (FOL): 2}

\begin{nbdef}{Provability $\Gamma \proves \varphi$ (page 2)}
Let $\Gamma$ be a set of formulas (the hypotheses). We write $\Gamma \proves \varphi$
(``$\Gamma$ proves $\varphi$'') if there is a finite sequence of formulas, each coming from
$\Gamma$, from the axioms, or from earlier formulas by an inference rule, ending at
$\varphi$.
\end{nbdef}

Recall that $\Gamma \entails \varphi$ (``$\Gamma$ entails $\varphi$'') means: in every
structure (interpretation) $M$ in which all formulas of $\Gamma$ are true, $\varphi$ is
true too. $\proves$ is about \emph{symbols and rules}; $\entails$ is about \emph{meaning and
truth}.

\begin{nbfact}{Soundness theorem (page 2)}
If $\Gamma \proves \varphi$, then $\Gamma \entails \varphi$. \quad (Everything provable is true.)
\end{nbfact}

\paragraph{Proof idea.} Take a proof $P$ of $\varphi$ from $\Gamma$ and a structure $M$ in
which $\Gamma$ is true, i.e.\ for all object assignments $\sigma$ (values for the free
variables) and all $\gamma \in \Gamma$, $M \entails \gamma[\sigma]$. Show by
\textbf{induction on the number of lines} of $P$ that every line $\psi$ is true in $M$
under all $\sigma$. Base cases: hypotheses (true by assumption) and axioms (true in every
structure). Modus Ponens preserves truth (truth table of $\limp$). For the
$\exists$-rule $\frac{\varphi \limp \psi}{\exists x\varphi \limp \psi}$ argue by
contrapositive:
\begin{enumerate}[nosep]
  \item suppose $M \nentails (\exists x\varphi \limp \psi)[\sigma]$; then
        $M \entails \exists x\varphi[\sigma]$ and $M \nentails \psi[\sigma]$;
  \item the first fact gives an object $m$ of $M$ with $M \entails \varphi[\sigma(m/x)]$
        ($\sigma$ changed so that $x$ denotes $m$);
  \item let $\sigma' = \sigma(m/x)$. Since $x$ is not free in $\psi$, changing the value of
        $x$ does not affect $\psi$, so $M \nentails \psi[\sigma']$;
  \item under $\sigma'$, $\varphi$ is true and $\psi$ false, so
        $M \nentails (\varphi \limp \psi)[\sigma']$: the premise of the rule is not true
        under all assignments. Contrapositive: if the premise is, so is the conclusion.
\end{enumerate}
(The notes write $M \entails \psi[\sigma(m/x)]$ in step 2: it should be $\varphi$.) Note
that we needed to \emph{change} the object assignment to make this work.

\begin{nbfact}{Completeness (G\"odel, 1929)}
\notinslides
The converse also holds: if $\Gamma \entails \varphi$ then $\Gamma \proves \varphi$. So
for FOL, ``provable'' and ``true in all models'' coincide. This does not contradict
undecidability: a proof always \emph{exists} for valid formulas, but there is no
algorithm that also tells you, in finite time, when \emph{no} proof exists.
\end{nbfact}

\begin{center}
\begin{tikzpicture}[font=\small, >=Stealth]
  \node[draw, rounded corners, thick, nbpurple, fill=nbpurple!6, align=center, minimum width=3.5cm, minimum height=1.1cm] (s) at (0,0) {$\Gamma \proves \varphi$\\\footnotesize syntax: there is a proof};
  \node[draw, rounded corners, thick, nbgreen, fill=nbgreen!6, align=center, minimum width=3.5cm, minimum height=1.1cm] (m) at (6.5,0) {$\Gamma \entails \varphi$\\\footnotesize semantics: true in every model};
  \draw[->, thick] (s.north east) to[bend left=15] node[above]{soundness} (m.north west);
  \draw[->, thick] (m.south west) to[bend left=15] node[below]{completeness} (s.south east);
\end{tikzpicture}
\captionof{figure}{Soundness: the rules never prove something false. Completeness: every
truth that follows has a proof.}
\end{center}

\subsubsection{Sequent calculus LK}
\slidesrc{Proof theory (FOL): 2--4}

A second style of proof system, invented by Gentzen: the \textbf{sequent calculus}. A
\textbf{sequent} has the form
\[ \Gamma \to \Delta \qquad (\Gamma, \Delta \text{ finite lists of formulas}) \]
and means: \emph{if all formulas in $\Gamma$ are true, then at least one formula in
$\Delta$ is true}. LK (from the German \emph{Logischer Kalk\"ul}) has the opposite
philosophy of Frege systems: \emph{one axiom, many rules}.

\paragraph{Variable conventions.} $a, b, c, \ldots$ are used only for \emph{free}
variables and $x, y, z, \ldots$ only for \emph{bound} variables. Consequence: a
sub-formula of a formula is not always a formula. In $\forall x\,(x \leq a)$, the part
$x \leq a$ is a \emph{semi-formula} ($x$ is not bound there); similarly $x + a$ is a
\emph{semi-term}. With these conventions any term $t$ is always substitutable for a free
variable $a$ in a formula $\varphi$.

\begin{nbfact}{LK (page 3)}
\textbf{Initial sequents}: $A \to A$, for any formula $A$. \\
\textbf{Rules}: all rules of the propositional sequent calculus, plus four quantifier rules:
\begin{center}
\begin{tabular}{cc}
\AxiomC{$\Gamma \to \Delta, A(b)$}
\LeftLabel{\small $\forall$-right}
\UnaryInfC{$\Gamma \to \Delta, \forall x A(x)$}
\DisplayProof
&
\AxiomC{$\Gamma \to \Delta, A(t)$}
\LeftLabel{\small $\exists$-right}
\UnaryInfC{$\Gamma \to \Delta, \exists x A(x)$}
\DisplayProof
\\[18pt]
\AxiomC{$A(b), \Gamma \to \Delta$}
\LeftLabel{\small $\exists$-left}
\UnaryInfC{$\exists x A(x), \Gamma \to \Delta$}
\DisplayProof
&
\AxiomC{$A(t), \Gamma \to \Delta$}
\LeftLabel{\small $\forall$-left}
\UnaryInfC{$\forall x A(x), \Gamma \to \Delta$}
\DisplayProof
\end{tabular}
\end{center}
In $\forall$-right and $\exists$-left the free variable $b$ is an \textbf{eigenvariable}:
it must not appear anywhere in the lower sequent (the substitution $x/b$ ``captures'' all
occurrences of $b$). In $\exists$-right and $\forall$-left, $t$ is \emph{any} term and may
occur in the lower sequent.
\end{nbfact}

\begin{nbintuition}[Eigenvariable = ``an arbitrary object we know nothing about'']
\begin{itemize}[nosep]
  \item \emph{$\forall$-right}: to prove ``for all $x$, $A(x)$'', prove $A(b)$ for a
        \emph{fresh} $b$ about which the hypotheses say nothing. Like in maths: ``let $b$
        be an arbitrary natural number\ldots''. If $b$ appeared in $\Gamma$ (e.g.\
        $\Gamma$ says $b = 0$), it would not be arbitrary.
  \item \emph{$\exists$-left}: if we know ``some $x$ has $A$'', we may give it a fresh name
        $b$, but we cannot assume it is any object we already know.
  \item \emph{$\exists$-right} and \emph{$\forall$-left}: to show ``something has $A$'',
        exhibit any witness $t$; to use ``everything has $A$'', instantiate with any $t$.
\end{itemize}
\end{nbintuition}

\subsubsection{LKe: adding equality}
\slidesrc{Proof theory (FOL): 4}

LK has no rules for $=$. \textbf{LKe} adds five initial sequents (for any terms
$s, t, u, s_i, t_i$):
\begin{itemize}[nosep]
  \item $\to s = s$
  \item $s = t \to t = s$
  \item $s = t,\ t = u \to s = u$
  \item $s_1 = t_1, \ldots, s_k = t_k \to f(s_1, \ldots, s_k) = f(t_1, \ldots, t_k)$ for any
        \emph{function symbol} $f$ of arity $k$
  \item $s_1 = t_1, \ldots, s_k = t_k,\ P(s_1, \ldots, s_k) \to P(t_1, \ldots, t_k)$ for
        any relation $P$ of arity $k$
\end{itemize}
(The notes write ``for any formula $f$'': $f$ is a function symbol. The last sequent is
printed with a second arrow; the version above is the intended one.)

\begin{nbexample}{proving symmetry of equality in LKe (page 4)}
Goal: $\to \forall x \forall y\,(x = y \limp y = x)$. Read the proof bottom-up as a plan,
top-down as a derivation:
\begin{center}
\AxiomC{$a = b \to b = a$}
\RightLabel{\small $\limp$-right}
\UnaryInfC{$\to a = b \limp b = a$}
\RightLabel{\small $\forall$-right ($b$ fresh)}
\UnaryInfC{$\to \forall y\,(a = y \limp y = a)$}
\RightLabel{\small $\forall$-right ($a$ fresh)}
\UnaryInfC{$\to \forall x \forall y\,(x = y \limp y = x)$}
\DisplayProof
\end{center}
\begin{enumerate}[nosep]
  \item Start from the LKe initial sequent $a = b \to b = a$ (symmetry for the terms $a, b$).
  \item $\limp$-right (a propositional rule: from $A, \Gamma \to \Delta, B$ infer
        $\Gamma \to \Delta, A \limp B$) moves $a = b$ to the right as an antecedent.
  \item $\forall$-right with eigenvariable $b$: $b$ does not occur in the lower sequent
        $\to \forall y\,(a = y \limp y = a)$, so the condition holds.
  \item $\forall$-right again with eigenvariable $a$.
\end{enumerate}
\end{nbexample}

\subsubsection{Natural deduction rules for quantifiers}
\notinslides

To connect with the natural deduction of Section~\ref{sec:natural-deduction-prop}, here are
the four quantifier rules in the same Intro/Elim style (same ideas as LK):
\begin{center}
\small
\renewcommand{\arraystretch}{1.3}
\begin{tabularx}{\linewidth}{@{}l X@{}}
\toprule
$\forall$\,Elim & from $\forall x\,\varphi(x)$ write $\varphi(t)$ for any term $t$ (instantiation) \\
$\forall$\,Intro & if you derived $\varphi(b)$ for a fresh $b$ appearing in no open
                   assumption, write $\forall x\,\varphi(x)$ \\
$\exists$\,Intro & from $\varphi(t)$ write $\exists x\,\varphi(x)$ \\
$\exists$\,Elim & from $\exists x\,\varphi(x)$, assume $\varphi(b)$ with $b$ fresh in a
                  subproof; if you derive $\psi$ (not containing $b$), write $\psi$ outside \\
\bottomrule
\end{tabularx}
\end{center}

\begin{nbexample}{the monitor argument, proved}
\begin{fitchproof}
\hypo{1}{\forall x\,(\mathit{Monitor}(x) \limp \mathit{Ready}(x))} \RPrem
\hypo{2}{\mathit{Monitor}(X12)} \RPrem
\have{3}{\mathit{Monitor}(X12) \limp \mathit{Ready}(X12)} \by{$\forall$\,Elim}{1}
\have{4}{\mathit{Ready}(X12)} \RImpE{3,2}
\end{fitchproof}
The argument of Section~\ref{sec:why-fol}, which propositional logic could not capture, becomes a two-step proof.
\end{nbexample}
